\chapter{Tasking and Concurrency}
\label{ch:tasking}

Concurrency in Ada is part of the language, not a bolted-on RTOS API: \emph{tasks}
for independent threads of control, \emph{protected objects} for race-free sharing,
interrupt handlers as protected procedures, and a monotonic clock for real-time
work. This chapter shows the practical mechanics on the dual-core ESP32-S3 --
creating tasks and interrupt handlers, periodic scheduling, waiting with timeouts,
the recurring concurrency patterns, watchdogs, and putting both cores to work. It
draws on the encapsulation features of Chapter~\ref{ch:encap} and the drivers of
the HAL chapters.

\section{Tasks and interrupt handlers}
\label{sec:tasks}\index{tasks}

Concurrency in Ada is in the language, not an RTOS API: you declare \emph{tasks}
for independent threads of control and \emph{interrupt handlers} as protected
procedures. This section shows the practical mechanics; the Interrupts chapter has
the full story on how attachment works underneath.

\subsection{Creating and starting a task}

You do not ``start'' an Ada task with a call -- you \emph{declare} it, and it
activates automatically. A single task is declared with its body; it begins
running when its master (the enclosing scope) reaches the code after the
declarations. For a \emph{library-level} task that master is the program itself,
so the binder activates it before \texttt{Main} even begins.

\cnote{there is no \texttt{xTaskCreate}, \texttt{pthread\_create}, or
\texttt{vTaskStartScheduler}: you declare the task and it activates itself, and
the Ada runtime \emph{is} the scheduler -- already running when \texttt{Main}
begins.}

\begin{lstlisting}
   task Blinker;                       -- declaration: this creates the task
   task body Blinker is
   begin                              -- runs concurrently once activated
      loop
         ESP32S3.GPIO.Toggle (2);
         delay until Clock + Milliseconds (250);
      end loop;
   end Blinker;
\end{lstlisting}

Add a \texttt{pragma Priority} (or the \texttt{Priority} aspect) to place it in
the priority order, and on this dual-core target pin it to a core with
\texttt{CPU => 1} or \texttt{CPU => 2}:

\begin{lstlisting}
   task Sampler with Priority => 5, CPU => 2;   -- runs on the second core
\end{lstlisting}

For a \emph{family} of identical workers use a task \emph{type} and declare
objects of it. One restriction matters here: under the \texttt{light-tasking} and
\texttt{embedded} (Jorvik) profiles a task may not be nested inside a subprogram
(\texttt{No\_Task\_Hierarchy}), so application tasks live at \emph{library level}
-- in a package -- not inside \texttt{Main}.

\subsection{Creating and attaching an interrupt handler}

An interrupt handler is a \emph{protected procedure} bound to an interrupt. The
protected object's \texttt{Interrupt\_Priority} is the hardware level it runs at
(and the ceiling that serialises it), and \texttt{pragma Attach\_Handler} wires
it up at elaboration -- no run-time call (\S\ref{ch:registers} and the Interrupts
chapter):

\begin{lstlisting}
   protected Handler with Interrupt_Priority => Device_L3_Priority is
   private
      procedure ISR;
      pragma Attach_Handler (ISR, Device_L3_0);   -- bound to the L3 device slot
   end Handler;
\end{lstlisting}

For GPIO pins the HAL wraps that boilerplate: \texttt{ESP32S3.GPIO.Interrupts}
owns the one GPIO interrupt source, routes it to the level-3 slot, and demuxes by
pin to per-pin callbacks. You just register a callback with a trigger:

\begin{lstlisting}
   Enable (Pin => 0, On => Falling_Edge, Action => On_Edge'Access);
\end{lstlisting}

A handler -- protected \texttt{ISR} or HAL callback -- runs in \emph{interrupt
context} at the slot's ceiling. Keep it tiny, and do not call a lower-ceiling
protected object or block. The idiom is to record the event (bump an
\texttt{Atomic}, or \texttt{Set} a \texttt{Suspension\_Object}) and let an
ordinary task do the real work.

\subsection{Worked example: counting falling edges on GPIO0}

Putting both together: a counter that increments on every falling edge of GPIO0,
with a task that reports it once a second. The interrupt callback is the only
writer of \texttt{Edges} (the level-3 ISR cannot re-enter itself), so an
\texttt{Atomic} \texttt{Natural} written there and read by the task needs no lock.
Everything is library-level, so it is legal on the \texttt{embedded} profile.

\begin{lstlisting}
with Ada.Real_Time;           use Ada.Real_Time;
with Ada.Text_IO;             use Ada.Text_IO;
with ESP32S3.GPIO;            use ESP32S3.GPIO;
with ESP32S3.GPIO.Interrupts; use ESP32S3.GPIO.Interrupts;

package body Edge_Demo is

   Edges : Natural := 0 with Atomic;     -- written by the ISR, read by the task

   --  Runs in interrupt context (level-3 ceiling): do the minimum.
   procedure On_Edge is
   begin
      Edges := Edges + 1;                -- sole writer -> the inc is safe
   end On_Edge;

   --  A normal task, activated automatically at start-up.
   task Reporter;
   task body Reporter is
   begin
      loop
         delay until Clock + Seconds (1);
         Put_Line ("GPIO0 falling edges:" & Natural'Image (Edges));
      end loop;
   end Reporter;

begin
   --  Package elaboration: configure the pin and attach the interrupt.
   Configure (0, Mode => Input, Pull => Pull_Up);          -- pulled-up input
   Enable (Pin => 0, On => Falling_Edge, Action => On_Edge'Access);
end Edge_Demo;
\end{lstlisting}

\paragraph{Reading it line by line.} For a reader new to Ada, here is every
construct in the example:

\begin{description}[style=nextline]
\item[\texttt{with Ada.Real\_Time;}] names a unit this code \emph{depends on} and
  makes its declarations available -- like an \texttt{\#include}, but semantic:
  the compiler records the dependency rather than pasting text. The four
  \texttt{with}s pull in timing (\texttt{Ada.Real\_Time}), console output
  (\texttt{Ada.Text\_IO}), the GPIO driver, and its interrupt child.
\item[\texttt{use Ada.Real\_Time;}] makes that unit's names usable \emph{without}
  the package prefix (see below). \texttt{with} grants access; \texttt{use} grants
  short names.
\item[\texttt{package body Edge\_Demo is}] begins the implementation of a package
  called \texttt{Edge\_Demo}. A package groups related declarations; everything up
  to \texttt{end Edge\_Demo;} belongs to it. The task it declares is legal under the Jorvik profile because it is a top-level
(\emph{library-level}) package.
\item[\texttt{Edges : Natural := 0 with Atomic;}] declares a variable. The colon
  separates the name from its type; \texttt{Natural} is the predefined type of
  integers $\geq 0$; \texttt{:= 0} sets its initial value (\texttt{:=} is
  \emph{assignment}; plain \texttt{=} is the equality \emph{test}). \texttt{with
  Atomic} is an aspect making each read or write a single indivisible access.
\item[\texttt{-- \ldots}] a comment: \texttt{--} runs to the end of the line.
\item[\texttt{procedure On\_Edge is \ldots\ begin \ldots\ end On\_Edge;}] defines a
  \emph{procedure} -- a subprogram that returns no value -- taking no
  parameters. The statements between \texttt{begin} and \texttt{end} run when it
  is called. Repeating the name after \texttt{end} is optional but conventional.
\item[\texttt{Edges := Edges + 1;}] an assignment statement: evaluate
  \texttt{Edges + 1} and store it back. Statements end with a semicolon.
\item[\texttt{task Reporter;}] declares a \emph{task} -- an independent thread of
  control. This line creates it; it will start on its own (no ``start'' call).
\item[\texttt{task body Reporter is \ldots\ loop \ldots\ end loop; \ldots}] the
  task's code. \texttt{loop \ldots\ end loop;} repeats indefinitely.
\item[\texttt{delay until Clock + Seconds (1);}] suspends the task until an
  \emph{absolute} instant: the current time (\texttt{Clock}) plus one second
  (\texttt{Seconds (1)} is a time span). \texttt{delay until} waits to a point in
  time; plain \texttt{delay X} would wait a duration.
\item[\texttt{Put\_Line ("\ldots" \& Natural'Image (Edges));}] prints a line.
  \texttt{\&} concatenates strings, and \texttt{Natural'Image (Edges)} converts
  the number to text. The apostrophe introduces an \emph{attribute} -- a built-in
  query on a type or object; \texttt{'Image} is ``the textual image of this
  value.''
\item[\texttt{begin \ldots\ end Edge\_Demo;}] a package body may end with an
  executable part. It runs \emph{once, automatically}, when the package is
  elaborated at start-up (before \texttt{Main}). Here it brings up the hardware.
\item[\texttt{Configure (0, Mode => Input, Pull => Pull\_Up);}] calls the GPIO
  driver. \texttt{0} is a positional argument (the pin); \texttt{Mode => Input}
  and \texttt{Pull => Pull\_Up} are \emph{named} associations
  (\texttt{name => value}), so the call documents which parameter is which.
\item[\texttt{Enable (\ldots, Action => On\_Edge'Access);}] registers the
  interrupt. \texttt{On\_Edge'Access} is the \texttt{'Access} attribute -- a
  reference to the \texttt{On\_Edge} procedure -- which the HAL stores and calls
  from the ISR on each falling edge.
\end{description}

\paragraph{Why \texttt{use} makes it readable.} A \texttt{use} clause makes a
package's names directly visible, so you write \texttt{Configure},
\texttt{Falling\_Edge} and \texttt{Put\_Line} instead of their full paths. Without
it, the same two setup lines would read:

\begin{lstlisting}
   ESP32S3.GPIO.Configure (0, Mode => ESP32S3.GPIO.Input,
                              Pull => ESP32S3.GPIO.Pull_Up);
   ESP32S3.GPIO.Interrupts.Enable
     (Pin => 0, On => ESP32S3.GPIO.Interrupts.Falling_Edge,
      Action => On_Edge'Access);
\end{lstlisting}

The package path repeats on nearly every name and buries the intent. With
\texttt{use} the code reads in the language of the problem -- ``configure pin 0
as a pulled-up input; enable a falling-edge interrupt'' -- while the \texttt{with}
clauses still record exactly where those names come from. The trade-off is that a
name's origin is less explicit, and two \texttt{use}d packages could in principle
declare the same name; in a small, focused unit like this the clarity gain wins,
and you can always write a single name in full where it aids understanding.

The application's \texttt{Main} need only \texttt{with Edge\_Demo;} and stay
alive; \texttt{Reporter} and the interrupt are running before \texttt{Main}'s
first statement. Each time GPIO0 goes high$\rightarrow$low (release of a pull-up
button, say) the ISR bumps \texttt{Edges}, and the reporter prints the running
total -- a complete interrupt-driven counter in a couple of dozen lines, with no
RTOS, no manual ISR vectoring, and no explicit task start.

\section{Real-time tasks: periodic and time-driven}
\label{sec:realtime}\index{real-time tasks}

For periodic work use \texttt{Ada.Real\_Time} -- a \emph{monotonic} clock that
never jumps -- not \texttt{Ada.Calendar}. The canonical cyclic task computes its
\emph{next} release as an absolute instant and \texttt{delay until}s it, so timing
error does not accumulate:

\begin{lstlisting}
   task body Sampler is
      Period : constant Time_Span := Milliseconds (10);
      Next   : Time := Clock;             -- first release = now
   begin
      loop
         Read_And_Process;
         Next := Next + Period;           -- the next absolute instant
         delay until Next;                -- no drift, however long the body took
      end loop;
   end Sampler;
\end{lstlisting}

Because \texttt{Next} advances by exactly \texttt{Period} regardless of how long
\texttt{Read\_And\_Process} ran, the task holds its average rate with no cumulative
jitter -- the crucial difference from a naive \texttt{delay Period}, which slips
by the body's run time every cycle. Two more tools are available on all profiles:

\begin{itemize}
  \item \texttt{Time\_Span} arithmetic (\texttt{Milliseconds},
        \texttt{Microseconds}, \texttt{To\_Duration}) for expressing intervals
        exactly;
  \item \texttt{Ada.Real\_Time.Timing\_Events} -- register a procedure to run at
        a future instant without dedicating a task, handy for a one-shot timeout.
\end{itemize}

\texttt{Ada.Calendar} exists too, for wall-clock dates, but its absolute form is
unsuited as a scheduling base (see the Limitations chapter) and on the bare
runtime it is anchored only to boot. Schedule with \texttt{Ada.Real\_Time}.

\section{Waiting for an event -- and timing out}
\label{sec:waiting}\index{waiting for an event}

Two questions come up constantly in firmware: how do I \emph{wait} for something
to happen -- a GPIO pin to change, a byte to arrive on the UART -- and how do I
\emph{give up} after a while? There are two ways to wait (poll, or block until an
interrupt wakes you), and the clean way to bound the wait with a timeout differs
across the three profiles because the leaner two forbid the \texttt{select}
statement.

\subsection{Polling (any profile)}

The simplest wait re-checks a condition in a loop, backing off with
\texttt{delay until} so other tasks run; a deadline turns it into a timeout. To
wait for a UART byte:

\begin{lstlisting}
with Ada.Real_Time; use Ada.Real_Time;
with ESP32S3.UART;  use ESP32S3.UART;

function Wait_For_Byte (S : Session; Timeout : Time_Span) return Boolean is
   Deadline : constant Time := Clock + Timeout;
begin
   while Available (S) = 0 loop
      if Clock >= Deadline then
         return False;                       -- gave up
      end if;
      delay until Clock + Milliseconds (1);  -- back off; let other tasks run
   end loop;
   return True;                              -- a byte is waiting
end Wait_For_Byte;
\end{lstlisting}

The same shape waits on a GPIO level -- replace \texttt{Available (S) = 0} with
\texttt{not ESP32S3.GPIO.Read (Pin)}. Polling works on \emph{every} profile and
needs only \texttt{Ada.Real\_Time}; the cost is the CPU spent re-checking.

Rather than hand-roll the deadline loop each time, the HAL bundles it as
\texttt{ESP32S3.Wait.Until\_True} (lock-free, all profiles): give it a condition
function and a timeout and it returns \texttt{True} on success or \texttt{False}
on timeout.

\begin{lstlisting}
with ESP32S3.Wait;
function Byte_Ready return Boolean is (Available (S) > 0);
...
if ESP32S3.Wait.Until_True (Byte_Ready'Access, Milliseconds (200)) then
   Read (S, Buf, N);     -- a byte arrived in time
end if;
\end{lstlisting}

\subsection{Blocking until an interrupt signals you (any profile)}

Better, when the wait may be long or power matters: sleep until the hardware wakes
you. Attach an interrupt handler that signals a \texttt{Suspension\_Object}, and
suspend the task on it:

\begin{lstlisting}
with Ada.Synchronous_Task_Control;  use Ada.Synchronous_Task_Control;
with ESP32S3.GPIO.Interrupts;

Edge_Seen : Suspension_Object;

procedure On_Edge is               -- runs in interrupt context: keep it tiny
begin
   Set_True (Edge_Seen);
end On_Edge;
...
ESP32S3.GPIO.Interrupts.Enable
  (Pin => 0, On => ESP32S3.GPIO.Interrupts.Falling_Edge, Action => On_Edge'Access);
...
Suspend_Until_True (Edge_Seen);    -- the task sleeps here until the edge fires
\end{lstlisting}

This burns no CPU while waiting. But \texttt{Suspend\_Until\_True} has no timeout
of its own -- and \emph{that} is where the profiles diverge.

\subsection{Timing out: three profiles, three approaches}

\textbf{Jorvik (light-tasking).} No \texttt{select} and no exception propagation,
so the timeout is a \emph{value} and the wait blocks on a protected entry whose
barrier opens on either the event or a deadline. The deadline is an
\texttt{Ada.Real\_Time.Timing\_Events} handler; the task genuinely sleeps, and on
return a flag says which fired:

\begin{lstlisting}
with Ada.Real_Time.Timing_Events; use Ada.Real_Time.Timing_Events;

protected Gate is
   procedure Signal;                            -- from the interrupt callback
   procedure Expire (E : in out Timing_Event);  -- the deadline handler
   entry Wait (Timed_Out : out Boolean);
private
   Occurred, Expired : Boolean := False;
end Gate;

protected body Gate is
   procedure Signal is begin Occurred := True; end Signal;
   procedure Expire (E : in out Timing_Event) is begin Expired := True; end Expire;
   entry Wait (Timed_Out : out Boolean) when Occurred or Expired is
   begin
      Timed_Out := not Occurred;
      Occurred := False;  Expired := False;     -- re-arm for next time
   end Wait;
end Gate;

--  Use: the interrupt callback calls Gate.Signal; the caller arms a deadline.
declare
   Deadline : Timing_Event;
   Late     : Boolean;
begin
   Set_Handler (Deadline, Clock + Milliseconds (500), Gate.Expire'Access);
   Gate.Wait (Late);          -- sleeps until the edge OR 500 ms, whichever first
   if Late then ... else ... end if;
end;
\end{lstlisting}

\textbf{embedded.} The same Jorvik primitives (still no \texttt{select}), but
exceptions now propagate -- so a timeout reads naturally as a \emph{catchable
exception}, giving linear code, with any controlled handles released on the way
out:

\begin{lstlisting}
Timeout_Error : exception;

procedure Wait_For_Byte (S : Session; Timeout : Time_Span) is
   Deadline : constant Time := Clock + Timeout;
begin
   while Available (S) = 0 loop
      if Clock >= Deadline then
         raise Timeout_Error with "no byte within the deadline";
      end if;
      delay until Clock + Milliseconds (1);
   end loop;
end Wait_For_Byte;
...
begin
   Wait_For_Byte (S, Milliseconds (200));
   Read (S, Buf, N);                 -- a byte is here
exception
   when Timeout_Error =>             -- handle the timeout in one place
      ...
end;
\end{lstlisting}

(You can raise the same exception from the \texttt{Gate.Wait} barrier above
instead of polling, if you would rather sleep than spin.)

\textbf{full.} Full Ada gives the timeout directly, as a \emph{timed entry call}:
the task blocks on the entry and a \texttt{delay} races it -- no helper deadline,
no polling. The protected object only needs the event half:

\begin{lstlisting}
protected Gate is
   procedure Signal;          -- from the interrupt callback
   entry Wait;
private
   Occurred : Boolean := False;
end Gate;

protected body Gate is
   procedure Signal is begin Occurred := True; end Signal;
   entry Wait when Occurred is begin Occurred := False; end Wait;
end Gate;

--  Use: whichever happens first wins.
select
   Gate.Wait;                              -- the event arrived
   ...
or
   delay until Clock + Milliseconds (500); -- the timeout
   ...
end select;
\end{lstlisting}

(Use \texttt{delay until} -- an \emph{absolute} time -- here: on this port a
\emph{relative} \texttt{delay} as a \texttt{select} trigger has a known gap, see
the Limitations chapter.) \texttt{full} also offers \texttt{select \dots\ then
abort} to put a deadline around an \emph{arbitrary} computation, not just an entry
call -- abandoning it if it overruns.

\textbf{Choosing.} Poll for a short, infrequent wait or the leanest code;
interrupt-signalled blocking is better when the wait may be long or power matters.
For the timeout: a status value on Jorvik, a catchable exception on embedded, and
the \texttt{select}\,/\,\texttt{delay} timed entry call on full.

\section{Choosing how to wait: a spectrum, by timescale}
\label{sec:wait-spectrum}\index{busy-wait}\index{polling!spectrum}\index{spin loop}

The previous section drew the choice as binary -- poll, or block on an interrupt.
That is the right coarse split, but a driver author actually picks among
\emph{four} points on a spectrum, and the deciding factor is almost always the
same question: \emph{how long is the wait expected to be?} Match the strategy to
the timescale and the rest follows.

\subsection{A tight spin -- shorter than a context switch}

When you are waiting on a hardware register bit that clears in a microsecond or
two -- a command latching, a DMA-fed shift register draining -- the right thing
is to spin with \emph{no} back-off at all:

\begin{lstlisting}
B.Regs.CMD.USR := True;            -- kick the SPI transfer
while B.Regs.CMD.USR loop null;    -- ... and wait the few us for it to finish
end loop;
\end{lstlisting}

This looks like the cardinal sin of embedded code -- a busy-loop burning the CPU
-- but here it is correct, and a \texttt{delay until} would be \emph{worse}. The
wait is a handful of microseconds; the kernel tick is a millisecond, and even a
bare context switch costs more cycles than the bit takes to clear. Sleeping would
turn a 2\,\textmu s wait into a 1\,ms one and add scheduler churn for nothing. The
SPI engine (\texttt{CMD.USR}, \texttt{CMD.UPDATE}), the I2S engine
(\texttt{TX\_UPDATE}), and the SHA peripheral (\texttt{BUSY.STATE}) all spin like
this, deliberately. The rule: \emph{if the wait is shorter than the cost of
yielding, do not yield.}

\subsection{A bounded spin -- a tight loop that cannot hang}

The pure \texttt{loop null} spin has one flaw: if the bit \emph{never} clears -- a
wedged peripheral, a chip that is not there -- it hangs the core forever. Where
that risk is real, the same spin gets a guard counter, so a fault becomes a
bounded failure instead of a deadlock:

\begin{lstlisting}
Guard : Natural := 0;
while not Done (C, Dir) and then Guard < 2_000_000 loop   -- GDMA completion
   Guard := Guard + 1;
end loop;
--  fall through whether the DMA finished or the cap tripped; the caller checks Done
\end{lstlisting}

The cap is a crude worst-case-execution-time bound: still no \texttt{delay} (the
expected wait is short), but liveness is guaranteed. GDMA's \texttt{Wait} and the
W5500's command-accept (\texttt{Issue}, capped at 2000 tries) use this shape.

\subsection{A paced poll -- when the wait is milliseconds}

Once the expected wait reaches the millisecond range -- a reset settling, a slow
status bit, a byte trickling in over a UART -- yielding pays off, and you back off
with \texttt{delay until} so other tasks run meanwhile. This is the polling loop
from the previous section, and \texttt{ESP32S3.Wait.Until\_True} packages it.
Datasheet-mandated settle delays (the ST7789's 120\,ms power-on, the W5500's reset
hold) are the same idea with a \emph{fixed} delay rather than a re-checked
condition.

\subsection{Blocking on an interrupt -- long or unbounded}

When the wait could be long, is unpredictable, or power matters -- a network
packet that may arrive in 1\,ms or 10\,s, a key-press, a sensor's data-ready line
-- polling wastes both CPU and energy. Sleep instead: an interrupt handler
\texttt{Set\_True}s a \texttt{Suspension\_Object} and the task
\texttt{Suspend\_Until\_True}s on it, burning nothing until the hardware wakes it
(Chapter~\ref{ch:suspension}). This is the top of the spectrum: the most machinery
(a handler, a wired interrupt), but the only choice that costs \emph{zero} while
waiting.

\subsection{One driver, the whole spectrum}

The W5500 networking stack is a useful specimen because it spans nearly all of
these at once. Accepting a chip command spins, bounded (\texttt{Issue}); a short
register read spins; and the blocking waits -- for an incoming connection, for
received data, for room in the transmit buffer -- funnel through a single hook:

\begin{lstlisting}
procedure Wait_Event (S : Socket) is
begin
   if Waiter /= null then Waiter (S.Index);     -- sleep on INTn (Suspend_Until_True)
   else delay until Clock + Milliseconds (1);   -- else a 1 ms paced poll
   end if;
end Wait_Event;
\end{lstlisting}

By default the stack \emph{polls} (the 1\,ms fallback); calling
\texttt{Interrupts.Enable} registers a waiter and the very same waits become
interrupt-blocking, with no change to the socket code. That default carries a
throughput cost worth knowing: a bulk transfer paced at one decision per
millisecond is far slower than one that wakes the instant the buffer drains -- the
difference the 1\,MiB FTP upload of Chapter~\ref{ch:networking} ran into. For
sustained data, enable the interrupts.

\subsection{The rule of thumb}

\begin{center}
\begin{tabular}{l l l}
\toprule
\textbf{Expected wait} & \textbf{Strategy} & \textbf{Example} \\ \midrule
under a few \textmu s          & tight spin, no back-off        & SPI \texttt{CMD.USR}, SHA \texttt{BUSY} \\
short, but may never end       & bounded spin (guard count)     & GDMA \texttt{Wait}, W5500 \texttt{Issue} \\
milliseconds                   & paced poll (\texttt{delay until}) & reset/settle, \texttt{Wait.Until\_True} \\
long / unbounded / power       & block on an interrupt          & W5500 \texttt{Wait\_Data}, sensor DRDY \\
\bottomrule
\end{tabular}
\end{center}

Two failure modes bracket the table. Yield on a sub-microsecond wait and you pay a
thousandfold latency penalty for nothing; spin -- even bounded -- on a
multi-second wait and you burn a core and a battery. Picking the row that matches
the timescale is most of getting it right.

\section{Concurrency patterns}
\label{sec:concurrency}\index{concurrency patterns}

A few protected-object shapes recur often enough to be worth naming. All work on
every tasking profile; the Jorvik restrictions (one entry per protected object,
one caller queued at a time) shape them but do not prevent them.

\textbf{Producer/consumer over a bounded buffer.} The classic decoupling: one
task (or an interrupt) produces items, another consumes them, with a protected
ring buffer between. The consumer's \texttt{Get} is an \emph{entry} whose barrier
holds it until something is available -- no polling.

\begin{lstlisting}
protected Queue is
   procedure Put (X : Item);          -- producer side (also callable from an ISR)
   entry     Get (X : out Item);      -- consumer side: blocks while empty
private
   Buf        : Item_Array (0 .. 15);
   Count      : Natural := 0;
   Head, Tail : Natural := 0;
end Queue;

protected body Queue is
   procedure Put (X : Item) is
   begin
      if Count < Buf'Length then      -- (full -> drop; or count an overflow)
         Buf (Tail) := X;
         Tail := (Tail + 1) mod Buf'Length;
         Count := Count + 1;
      end if;
   end Put;

   entry Get (X : out Item) when Count > 0 is
   begin
      X := Buf (Head);
      Head := (Head + 1) mod Buf'Length;
      Count := Count - 1;
   end Get;
end Queue;
\end{lstlisting}

The consumer task just loops \texttt{Queue.Get (X)}; it sleeps on the barrier when
the queue is empty and wakes the instant \texttt{Put} adds an item. (Under Jorvik,
exactly one consumer task may queue on \texttt{Get} at a time -- which is the
usual one-producer/one-consumer arrangement.)

\textbf{A signal / mailbox.} The degenerate one-slot case is a handoff: an
interrupt deposits a value and a task collects it. It is the
``protected entry'' half of the waiting pattern (\S\ref{sec:waiting}) carrying a
payload -- a \texttt{procedure Post (V)} from the ISR and an
\texttt{entry Take (V : out \ldots)} barriered on ``something posted.''

\textbf{DMA double-buffering (ping-pong).} For a continuous stream -- audio in
over I2S, samples to a DAC -- you cannot stop to process between transfers. Use
\emph{two} buffers: while the GDMA fills (or drains) buffer~A, the CPU processes
buffer~B; when the transfer completes you swap, so the peripheral is never idle.
The completion interrupt signals the processing task (a suspension object or a
queue of ``buffer ready'' indices) and points the next transfer at the other
buffer. The protected object owns the ``which buffer is the hardware's'' state;
the task owns the other. This keeps a real-time stream gap-free with bounded,
analysable hand-off.

\begin{figure}[ht]
\centering
\begin{tikzpicture}[node distance=8mm and 30mm]
\node[blk, fill=green!9, text width=3.0cm]  (a1) {Buffer A -- DMA fills};
\node[blk, below=of a1, text width=3.0cm]   (b1) {Buffer B -- CPU works};
\node[blk, right=of a1, fill=green!9, text width=3.0cm] (b2) {Buffer B -- DMA fills};
\node[blk, below=of b2, text width=3.0cm]   (a2) {Buffer A -- CPU works};
\draw[flow] (a1.east) -- (a2.west);     % A: DMA -> CPU
\draw[flow] (b1.east) -- (b2.west);     % B: CPU -> DMA
\node[font=\scriptsize, fill=white, inner sep=1pt] at ($(b1)!0.5!(b2)$) {swap on DMA-done IRQ};
\node[font=\scriptsize] at ($(a1.north)!0.5!(b1.south)-(2.4cm,0)$) {phase 1};
\node[font=\scriptsize] at ($(b2.north)!0.5!(a2.south)+(2.4cm,0)$) {phase 2};
\end{tikzpicture}
\caption{DMA double-buffering: the peripheral always has one buffer while the CPU
works the other; the completion interrupt swaps their roles (the crossing), so the
stream never stalls.}
\label{fig:pingpong}
\end{figure}

\textbf{Why protected objects, not flags.} Each of these could be hacked with a
shared \texttt{Volatile} flag, but the protected object gives you the barrier (no
polling), mutual exclusion (no torn reads across cores), and a priority ceiling
(bounded blocking) for free -- the points the Encapsulation chapter
(Chapter~\ref{ch:encap}) makes about \texttt{protected}.

\section{Watchdogs and recovery}
\label{sec:watchdogs}\index{watchdog}

A watchdog is a hardware timer that resets the chip unless the software keeps
``feeding'' it -- the last line of defence against a hang that no exception or
check can catch (a wild loop, a deadlock, a stuck peripheral). The ESP32-S3 has
several: the two timer-group main watchdogs (TG0/TG1 MWDT), the RTC watchdog
(RWDT), and the RTC \emph{super}-watchdog (a hardware backstop).

\textbf{The bare boot leaves them off.} The ROM and the second-stage bootloader
arm watchdogs during boot, but this runtime's startup \emph{disables} them (so a
long elaboration or a debugger pause is not cut short). So by default a bare-metal
app runs with \emph{no} active watchdog -- a hang will sit forever rather than
recover. That is fine while developing (and is why a wedged example does not
reset on its own), but a deployed device usually wants protection.

\textbf{Arming the RTC watchdog.} To get auto-recovery, enable the RWDT with a
timeout and feed it from a place that only runs while the system is healthy. The
registers are key-locked; the sequence is unlock, configure, feed, re-lock:

\begin{lstlisting}
--  RTC watchdog registers (overlaid as plain words).
Key     : constant := 16#50D8_3AA1#;
Config0 : Unsigned_32 with Volatile, Address => To_Address (16#6000_8098#), Import;
Config1 : Unsigned_32 with Volatile, Address => To_Address (16#6000_809C#), Import;
Feed    : Unsigned_32 with Volatile, Address => To_Address (16#6000_80AC#), Import;
Wprot   : Unsigned_32 with Volatile, Address => To_Address (16#6000_80B0#), Import;

procedure Arm (Timeout_Cycles : Unsigned_32) is   -- RTC-slow-clk (~63 kHz) cycles
begin
   Wprot   := Key;                       -- unlock
   Config1 := Timeout_Cycles;
   Config0 := 2**31                      -- WDT_EN
            or (3 * 2**28)               -- stage 0 action = reset the main system
            or (7 * 2**16) or (7 * 2**13);  -- reset-pulse lengths
   Feed    := 2**31;
   Wprot   := 0;                         -- re-lock
end Arm;

procedure Pet is begin Wprot := Key; Feed := 2**31; Wprot := 0; end Pet;
\end{lstlisting}

\textbf{Feed on progress, not on a timer.} The robust idiom is to pet the dog from
points that only execute when work is actually advancing -- the main loop body, a
completed transaction -- not from a periodic task that would keep petting even if
the real work has stalled. The ACATS test harness in this project does exactly
this: it arms the RWDT (\,$\approx$19\,s) and feeds it on every progress line, so a
test that deadlocks stops feeding and is reset, while a slow-but-progressing test
survives.

\textbf{The super-watchdog} is a separate hardware backstop that a deep-sleep wake
can leave armed; \texttt{ESP32S3.RTC.Disable\_Super\_Watchdog} turns it off when an
app means to stay awake (the RTC driver also re-feeds it automatically after a
deep-sleep wake). Between an armed RWDT for hang-recovery, the precise
stack-overflow watchpoint (the Debugging chapter), and a last-chance handler for
the unrecoverable case, a deployed app has a layered safety net.

\section{Dual-core (SMP) programming}
\label{sec:smp}\index{SMP}

The ESP32-S3 has two cores and the runtime schedules across both. Two ordinary-Ada
facilities put them to use.

\textbf{Pin a task to a core} with the \texttt{CPU} aspect (\texttt{CPU => 1} is
the first core, \texttt{2} the second); omit it and the task may run on either:

\begin{lstlisting}
   task Beat0 with Priority => Priority'Last - 1, CPU => 1;
   task Beat1 with Priority => Priority'Last - 1, CPU => 2;   -- the other core
\end{lstlisting}

\textbf{Share data with a protected object} -- it works \emph{across} cores, not
just between tasks on one core. A protected \emph{entry} can even hand a caller
from one core to the other: in the \texttt{esp32s3\_smp} example a producer on
core 1 blocks in \texttt{entry Get when Full}, and serving the entry wakes it
through the inter-core interrupt. The locking is the same priority-ceiling monitor
(Chapter~\ref{ch:encap}); the runtime adds the cross-core wake-up underneath.

\begin{lstlisting}
   protected Mailbox is
      procedure Put (V : Integer);
      entry     Get (V : out Integer);   -- blocks until Full; may serve cross-core
   private
      Item : Integer;
      Full : Boolean := False;
   end Mailbox;
\end{lstlisting}

Two cautions specific to this target: time comes from a \emph{shared} hardware
timer (the SYSTIMER), not each core's own cycle counter, so \texttt{Clock} agrees
across cores; and the USB-serial-JTAG console is not robust for two tasks printing
at once -- funnel diagnostics through one task. With those in mind, an SMP
application is just ordinary tasks, pinned where you want them, sharing through
protected objects.

